\documentclass[11pt]{article}
\usepackage{amsmath,amssymb}
\newcommand{\E}{\mathbb E}
\newcommand{\Cov}{\operatorname{Cov}}
\newcommand{\Var}{\operatorname{Var}}
\title{Literal support multiplicities and stopped elimination covariance}\author{}\date{}
\begin{document}\maketitle
\paragraph{Analysis of the linear support count in the local ordinary covariance calculation and the inherited effective score calculation.}
A four-dimensional bond belongs to six plaquettes. Its force reads exactly its own bond
and the other eighteen bonds in those plaquettes, on an embedded patch. Let $D(b)$ denote
this nineteen-bond determining set. A source row $c$ varies under a free link $e$ precisely
when $e\in D(c)$; incidence symmetry makes this equivalent to $c\in D(e)$.
For a two-link block, the varying row set is $D(e_-)\cup D(e_+)$, containing thirty-seven
bonds in the geometry used here. With the originally required disjoint determining
neighborhoods, these row sets are disjoint, so for a literal test
\[
 W(f)^2=\sum_j\sum_{c\in\mathcal R_j}|f_c|^2
 \le\sum_{c:f_c\ne0}|f_c|^2.
\]
Without disjointness the correct bound includes the actual maximal row multiplicity;
bounded lattice incidence alone bounds it by nineteen for a packing with distinct free links.
The disjoint packing used in the proof has no hidden volume-dependent multiplicity.

For the first-layer conditional mean of a single source row $c$, its retained carrier is
contained in
\[
 D_R(c)=(D(c)\setminus E)\cup
       \bigcup_{e\in E\cap D(c)}(D(e)\setminus\{e\}).
\]
Every new dependence comes from the actual one-link conditional denominator or numerator.
Thus $|D_R(c)|\le19+19\cdot18=361$, a crude fixed bound. Consequently each source-row
weight occurs in the summed derivative estimate of the inherited effective score calculation at most $361$ times per fixed
color count. This proves the claimed linear support dependence independently of how
many source rows or blocks are present. It does not assert an optimal geometric constant.

The fixed pointed completion contributes its literal finite number of transformed tests.
Applying the estimate before their finite average gives
$W^2\le C_{\rm point}\sum_{\gamma,c}|f_c^\gamma|^2$ with every shifted row retained.
If $Y_f$ is the original completed union of support plaquettes and $N_f=\#Y_f$, each
nonzero bond test is incident to a plaquette in that union; each plaquette has four bonds.
Therefore the relevant number of nonzero completed bond weights is at most
$4C_{\rm point}N_f$, and
$W^2\le C\|f\|_\infty^2N_f$. This is the precise convention under which the support-count
claim holds. Defining $N_f$ as only an uncompleted continuum support, or counting a
single representative before its offsets, would not justify the displayed inequality.
The original $a^4$ smearing is already in $\sigma=g^{-2}a^{-2}$; it is not inserted again
when converting this bond count to physical volume.

\paragraph{Exact reverse-martingale decomposition for true coordinate elimination.}
Let $R_0\supset R_1\supset\cdots\supset R_m$ be actual retained coordinate sets,
$\mathcal F_k=\sigma(U_b:b\in R_k)$, and $G_k=\E(G\mid\mathcal F_k)$.
Put $D_kG=G_{k-1}-G_k$. For every pair of square-integrable real observables,
\[
 \Cov(G,H)=\sum_{k=1}^m\E[(D_kG)(D_kH)]+\Cov(G_m,H_m),
\]
with $G_0=G$ when $R_0$ contains all coordinates. The increments are orthogonal because
$\E(D_kG\mid\mathcal F_k)=0$ and later increments are $\mathcal F_k$-measurable.
Each summand is exactly
\[
 \E\Cov(G_{k-1},H_{k-1}\mid\mathcal F_k).
\]
The first separated-link summand is quantitatively bounded by the local ordinary covariance calculation. The remainder
$\Cov(G_1,H_1)$ is exactly the inherited mean covariance, with the effective density and
score of the inherited effective score calculation. There is no sign cancellation available for the equal-label remainder:
all martingale variance contributions are nonnegative.

\paragraph{Separated supports.}
In the first layer, if the actual varying blocks of $G(f)$ and $G(h)$ are disjoint,
conditional independence gives
$\Cov(G(f),G(h)\mid\mathcal F_1)=0$ pointwise. Hence
\[
 \Cov_\mu(G(f),G(h))=\Cov_{\nu_1}(\overline G(f),\overline G(h)).
\]
This is an exact localization of the unresolved separated correlation to the inherited
boundary fields, not a decay estimate. Their carriers are the finite enlarged sets $D_R(c)$
above, so the separation condition is a checkable finite incidence statement with original
test offsets. In further layers one may stop before enlarged carriers meet and retain
all remaining covariance in the stopped term. No decay factor is produced merely by the
absence of an early overlap.

\paragraph{What persists, and what does not, at later layers.}
Every coordinate marginal still has the exact score and directional moment estimates in
the inherited effective score calculation. Moreover, varying one retained link changes the original action by at most a fixed
constant, so ratios of its marginal densities are bounded by $e^{C\beta}$, independently
of how many coordinates were eliminated. Thus a one-link conditional Haar comparison
provides a gap at least $c e^{-C\beta}$ at every coordinate marginal.
This is a real bound for the actual marginal, but it retains the exponential loss. It is
not a uniform tensorization constant for a layer of interacting retained variables.
After the first elimination, a term $-\log z_e$ can couple several remaining links.
Eliminating another independent set in the ORIGINAL plaquette graph need not give a
product law for this effective interaction. One must use its actual enlarged interaction
graph, carrying functions and conditional normalizers. These neighborhoods can grow.

Even if successive independent layers are selected in their actual graphs, the integrated
observables $G_k$ are conditional means, not fixed finite Wilson-word force squares.
The resampling proof in the local ordinary covariance calculation therefore cannot be copied to $G_k$ without estimates for
its actual response. The inherited effective score calculation supplies one first-layer derivative bound; converting it to a
variance via only the available gap costs $e^{C\beta}$. The score moment bounds control
directional averages but do not prove a lower Hessian bound or spatial summability.
Thus the martingale identity cleanly identifies the next quantitative estimate: bound the
conditional covariance of the first inherited means, under their genuine effective law,
with the required grade and physical mesh, while retaining the stopped boundary covariance.
The identity itself does not close MT36 or give a global/separated continuum estimate.
\end{document}
